\intro

Злокачественные новообразования являются одной из наиболее распространённых
причин смертности в мире \cite{major_death_causes}. Этот факт стимулирует
постоянный поиск новых методов лечения, среди которых выделяется иммунотерапия
\cite{immunotherapy_overview},
которая воздействует не на саму опухоль, а активирует для борьбы с ней
собственную иммунную систему пациента.

Однако, несмотря на свой потенциал, эффективность иммунотерапии может
существенно различаться у разных пациентов \cite{immunotherapy_efficacy}: начального
состояния иммунной системы пациента, характеристик опухоли и выбранной
терапевтической стратегии.
Для формализации таких сложных взаимосвязей и анализа динамики системы в работе
используется метод математического моделирования на основе системы обыкновенных
дифференциальных уравнений.

Выбор системы обыкновенных дифференциальных уравнений (ОДУ) в качестве механизма
описания взаимосвязей разных процессов, влияющих на опухоль, обусловлен их
способностью эффективно описывать среднюю динамику больших
клеточных популяций, сохраняя при этом относительную вычислительную простоту. В
отличие от стохастических или агент-ориентированных моделей, которые отслеживают
поведение отдельных клеток, модели на основе системы ОДУ позволяют быстро
проводить широкий параметрический анализ и исследовать общие закономерности
системы. Такой анализ включает исследование чувствительности
модели к изменению её ключевых параметров: скорость гибели иммунных
клеток, скорость пролиферации опухолевых клеток, константа скорости их
уничтожения или , что помогает выявить ключевые факторы, оказывающие наибольшее влияние на
эффективность терапии.

Цель данной работы заключается в разработке программного обеспечения для
моделирования динамики опухолевых процессов, которое основано на математической
модели в виде системы дифференциальных уравнений и описывает динамику популяций
описанных в нём клеток во времени, а также проведение на его основе набора
вычислительных экспериментов для установления взаимосвязи между различными
параметрами модели и динамикой роста опухоли.

Для достижения поставленной цели были определены следующие задачи:
\begin{enumerate}
  \item Описать математическую модель динамики роста опухоли с учетом иммунотерапии на
    основе системы дифференциальных уравнений.
  \item Описать и обосновать выбор численных методов, используемых для решения
    данной системы уравнений.
  \item Спроектировать и реализовать программное обеспечение с графическим
    пользовательским интерфейсом для моделирования динамики роста опухолей.
  \item Провести серию вычислительных экспериментов для установления влияния
    различных входных параметров модели на динамику роста опухоли.
  \item Провести сравнение полученных результатов численных экспериментов с данными
    клинических наблюдений.
\end{enumerate}
